% Linux Kniha kouzel, šablona formátu LaTeX A5
% Copyright (c) 2019 Singularis <singularis@volny.cz>
%
% Toto dílo je dílem svobodné kultury; můžete ho šířit a modifikovat pod
% podmínkami licence Creative Commons Attribution-ShareAlike 4.0 International
% vydané neziskovou organizací Creative Commons. Text licence je přiložený
% k tomuto projektu nebo ho můžete najít na webové adrese:
%
% https://creativecommons.org/licenses/by-sa/4.0/
%
% This program is free software; you can redistribute it and/or modify
% it under the terms of the GNU General Public License as published by
% the Free Software Foundation; either version 2 of the License, or
% (at your option) any later version.
%
% This program is distributed in the hope that it will be useful,
% but WITHOUT ANY WARRANTY; without even the implied warranty of
% MERCHANTABILITY or FITNESS FOR A PARTICULAR PURPOSE. See the
% GNU General Public License for more details.
%
% You should have received a copy of the GNU General Public License
% along with this program; if not, write to the Free Software
% Foundation, Inc., 51 Franklin St, Fifth Floor, Boston, MA  02110-1301  USA
%
{{ZAČÁTEK KNIHY}}
\documentclass[10pt,final]{book}
% 1. skupina balíčků:
\usepackage[utf8]{inputenc}
\usepackage[czech]{babel}
%
{{POKUD JE FORMÁT pdf-a4}}
\usepackage[pdftex,a4paper,inner=1.5cm,outer=1cm,top=1.5cm,bottom=1.5cm]{geometry} % + showframe pro ladění
{{KONEC POKUD}}
{{POKUD JE FORMÁT pdf-a5}}
\usepackage[pdftex,a5paper,inner=1.5cm,outer=1cm,top=1.5cm,bottom=1.5cm]{geometry} % + showframe pro ladění
{{KONEC POKUD}}

%\usepackage[B1,LY1,T1]{fontenc}
\usepackage[T1]{fontenc}

% 2. skupina balíčků:
\usepackage{textcomp}       % musí být načtený přednostně; vyžadovaný balíčkem {gensymb}
\usepackage[pdftex]{xcolor} % barevné písmo

% 3. skupina balíčků:
\usepackage[arrows]{boisik} % jen kvůli znaku $\carriagereturn$ (pokud existuje stabilní náhrada, bylo by lepší ji použít)
\usepackage{fancyhdr}       % formátování záhlaví a zápatí
\usepackage[bottom]{footmisc}% aby se poznámky pod čarou správně umístily
\usepackage{gensymb}        % symboly jako např. \degree
\usepackage[pdftex]{graphicx}% vkládání obrázků
\usepackage{ifthen}         % příkaz \ifthenelse a související konstrukce
\usepackage{titlesec}       % formátování nadpisů kapitol
\usepackage{verbatim}       % zejm. prostředí {comment}, ale také {verbatim}

% {fancyhdr}:
\fancypagestyle{normalni}{%
    \fancyhead{}\fancyfoot{}%
    \fancyfoot[LE,RO]{\thepage}%
    \renewcommand*{\headrulewidth}{0pt}%
    \renewcommand*{\footrulewidth}{0pt}%
}

% {titlesec}:
\titleformat{\chapter}[hang]{\normalfont\LARGE\bfseries}{\makebox[35pt][l]{\thechapter}}{5pt}{}\relax
\titlespacing*{\chapter}{0pt}{-1ex}{1ex plus1ex minus0.5ex}\relax

\titleformat{\section}[hang]{\normalfont\Large\bfseries}{\makebox[35pt][l]{\thesection}}{5pt}{}\relax
\titlespacing*{\section}{0pt}{1ex plus2ex minus0.5ex}{0.5ex plus0.25ex minus0.5ex}\relax

\titleformat{\subsection}[hang]{\normalfont\large}{}{0pt}{\filcenter\arabic{subsection}~}\relax
\titlespacing*{\subsection}{0pt}{1ex plus4ex minus0pt}{4pt}\relax

\assignpagestyle{\chapter}{normalni}

% {xcolor}:
\definecolor{seda}{gray}{0.5}

\newenvironment*{blok}{}{}
% \begin{priklad}{#2 = Záhlaví příkladu}{#3 = \footnote{}\footnote{}...}
%   \radekprikladu{text-řádku}
%   \radekprikladu{text-řádku}
% \end{priklad}
\newenvironment*{priklad}[3][]{%
    \newcommand*{\radekprikladu}[1]{%
        \nepoprve{{%
            \fontfamily{bsk}\fontseries{m}\fontshape{n}\selectfont%
            \textcolor{seda}{\ensuremath{\carriagereturn}}%
        }}%
        \item{##1}%
    }%
    \newcommand*{\akceprikladu}[1]{%
        \nepoprve{{%
            \fontfamily{bsk}\fontseries{m}\fontshape{n}\selectfont%
            \textcolor{seda}{\ensuremath{\carriagereturn}}%
        }}%
        \item\ensuremath{\Rightarrow}~{\rmfamily{}##1}%
    }%
    \begin{samepage}
    \fontfamily{qcr}\fontseries{m}\fontshape{n}\selectfont\raggedright%
    \ifthenelse{\equal{#2}{}}{}{\begin{list}{\makebox[0pt][l]{\#}\makebox[0pt][l]{\rule[2ex]{5cm}{0.5pt}}}{%
            \setlength{\leftmargin}{0.75em}%
            \setlength{\labelwidth}{0pt}%
            \setlength{\labelsep}{\leftmargin}%
        }\raggedright\item\relax%
        {#2}#3%
    \end{list}\vspace{-3ex}}%
    \budepoprve%
    \begin{list}{}{%
        \setlength{\itemsep}{0pt}%
        \setlength{\parsep}{0pt}%
        \setlength{\topsep}{1ex plus0.5ex minus0.5ex}%
}}{%
    \end{list}%
    \end{samepage}%
    \vspace{2mm plus2mm minus2mm}%
}

% qcr = TeX Gyre Cursor

\newcommand*{\kapitola}[2][]{\chapter{#2}}
\newcommand*{\sekce}[2][]{\section{#2}}
\newcommand*{\podsekce}[2][]{\subsection{#2}}

\newcommand*{\budepoprve}{%
    \renewcommand{\nepoprve}[1]{%
        \renewcommand{\nepoprve}[1]{####1}%
    }%
}
\newcommand*{\dopln}[1]{{\fontfamily{qcr}\fontseries{m}\fontshape{it}\selectfont\underline{{#1}}}}%
\newcommand*{\moznyzlom}{{\fontseries{m}\fontshape{n}\selectfont\textcolor{seda}{\-}}}%
\newcommand{\nepoprve}[1]{#1}
\newcommand*{\tritecky}{{\fontfamily{cmss}\fontseries{m}\fontshape{n}\selectfont\normalsize\mbox{...}}}%
% {} \big \Big \bigg \Bigg
\newcommand*{\volitelnyzacatek}{\textcolor{seda}{\ensuremath{\big[}}}%
\newcommand*{\volitelnykonec}{\textcolor{seda}{\ensuremath{\big]}}}%

\renewcommand*{\rmdefault}{cmr}

\begin{document}%
%
% TITULNÍ STRANA
%
\pagestyle{empty}%
\vspace*{0.05\textheight}%
\begin{center}\fontfamily{cmss}\selectfont\bfseries\Huge\scalebox{1.5}{LINUX}\\[0.01\textheight]\Large Kniha kouzel\end{center}%
\vspace{1cm}%
\begin{blok}%
    \setlength{\unitlength}{0.5\textwidth}%
    \begin{picture}(2,1.4)%
        \put(0,0){\includegraphics[width=0.5\textwidth]{../pdf-spolecne/_obrazky/logo-knihy-velke.png}}%
        \put(1,1.4){\parbox[t][0.7\textwidth][c]{0.5\textwidth}{\centering\itshape%
            Praktická sbírka řešených příkladů\\[0.02\textwidth]%
            Veškerá moc příkazové řádky/příkazového řádku přehledně,
            pro~začátečníky i pokročilé}}%
    \end{picture}%
\end{blok}%
\par\vfill%
\begin{center}\textcopyright\,2019 Singularis a ostatní přispěvatelé\end{center}%
\clearpage%
%
% DRUHÁ STRANA
%
\pagestyle{empty}%
{\noindent\Huge Linux --- Kniha kouzel\par}%
\vspace{2ex}\noindent%
\textcopyright~2019 Singularis a ostatní přispěvatelé%

\vspace{2ex}\noindent%
Toto dílo \quotedblbase Linux: Kniha kouzel\textquotedblleft{} podléhá licenci
Creative Commons Attribution-ShareAlike 4.0 International. Úplný text licence
je přiložen. Pro jeho zobrazení na počítači navštivte:
\begin{center}\fontfamily{qcr}\selectfont%
http://creativecommons.org/licenses/by-sa/4.0/
\end{center}

\noindent\textbf{Použitá díla jiných autorů:}
\begin{itemize}%
\item obrazky/ve-vystavbe.png \textcopyright~2010, Sarang (volné dílo podle německých zákonů)
\item tučňák Tux \textcopyright~Larry Ewing; \textcopyright~Garrett LaSage (licence CC0, upraveno)
\item ReText 7.0.1 (snímek obrazovky) \textcopyright~2011-2017 Dmitry Shachnev <mitya57@gmail.com> a přispěvatelé (licence GPL3), zobrazený text: \textcopyright~2019 Singularis (licence Creative Commons Attribution-ShareAlike 4.0 International; kompatibilní s GPL3)
\end{itemize}%
\vfill%
\clearpage%
%
% Nastavit styl záhlaví a zápatí
\pagestyle{normalni}%
% Vysázet obsah
\renewcommand*{\contentsname}{OBSAH}%
\tableofcontents%
%
% Nastavit znak pro „dělení slov“ v kódu na zpětné lomítko
{%
    \fontfamily{qcr}%
    \fontseries{m}\selectfont\hyphenchar\font=`\\%
    \fontseries{b}\selectfont\hyphenchar\font=`\\%
    \fontseries{bx}\selectfont\hyphenchar\font=`\\%
    \fontshape{it}\fontseries{m}\selectfont\hyphenchar\font=`\\%
    \fontseries{b}\selectfont\hyphenchar\font=`\\%
    \fontseries{bx}\selectfont\hyphenchar\font=`\\%
}%
%
\clearpage%
\raggedbottom%
{{ZAČÁTEK KAPITOLY}}
{{TĚLO KAPITOLY}}
{{KONEC KAPITOLY}}
\end{document}
{{KONEC KNIHY}}
